\subsection{A bridging model for parallel computation.}
The article describes the model of the bulk synchronous parallel computer (BPSC), which is meant to be a briding model decoupling the hardware and software aspects of parallel computation. The BPSC model is described as consisting of three components: a set of processing and memory units, a router to deliver messages between those units and a synchronization mechanism to coordinate sequences of computation steps (supersteps) globally. The focus of the article is put on dealing with the distribution of memory contents needed in the local computation steps, as well as the networking.\\
The article illustrates, how an efficient usage of the different components could be achieved with regards to parallelism, for example by proposing hashing functions to deal with questions of data distribution or starting a much higher number of processes then processors are available to create parallel slackness. It thus provides a hardware-agnostic model, which then would enable to reason about the efficiency of parallel algorithms without needing to take into account the specific details of the available architecture.\\
In our opinion, the BSPC model gives useful abstraction ideas, many of which are still relevant today. The simplicity of the undelying idea provides a good basis to reason about the basic challenges of parallel computation, e.g. data distribution and communication. However, it is limited in capturing the complexity of modern architectures, especially with regards to memory hierarchies and communication patterns.